\documentclass[11pt,a4paper]{article}
\usepackage[utf8]{inputenc}
\usepackage[T1]{fontenc}
\usepackage[french]{babel}
\usepackage{geometry}
\usepackage{hyperref}
\usepackage{array}
\usepackage{enumitem}
\geometry{margin=2.5cm}

\title{\textbf{Rapport -- Mini-Projet Compilation}\\
Conception et réalisation d’un compilateur pour un langage personnalisé}
\author{Réalisé par : Gharbi Eya, Berraies Mouhamed Khalil, Ben Farhat Saoussen, Machergui Nada\\
Classe : 1 ING 4}
\date{}

\begin{document}
\maketitle

\section*{1. Présentation du projet et DSL choisi}
Le projet consiste à construire un compilateur pour un DSL appelé \textbf{BotScript}.  
Ce langage contrôle un robot virtuel pour dessiner sur un canvas.  
Le traitement suit la chaîne classique : \textbf{lexeur → parseur → AST → analyse sémantique → interprétation}.

\section*{2. Choix technologiques : outils et langages}
\begin{itemize}[leftmargin=1.2cm]
  \item \textbf{Langage principal :} TypeScript.
  \item \textbf{Interface :} React + Vite (IDE web).
  \item \textbf{Analyse lexicale et syntaxique :} Flex/Bison.
  \item \textbf{Exécution :} interprétation directe.
  \item \textbf{Compilation WASM :} Emscripten pour compiler Flex/Bison en WebAssembly.
\end{itemize}

\section*{3. Structure des fichiers (fichiers importants)}
\begin{itemize}[leftmargin=1.2cm]
  \item \texttt{src/compiler.ts} : lexeur, parseur, AST, analyse sémantique, interpréteur.
  \item \texttt{src/App.tsx} : interface web.
  \item \texttt{src/specs/botscript.l} : règles Flex (lexer).
  \item \texttt{src/specs/botscript.y} : grammaire Bison + actions AST.
  \item \texttt{public/wasm/} : fichiers WebAssembly générés.
  \item \texttt{src/tests/} : tests unitaires et end-to-end.
\end{itemize}

\section*{4. Structure de la grammaire et de l’AST}

\subsection*{4.1 Grammaire (résumé)}
\begin{itemize}[leftmargin=1.2cm]
  \item \textbf{Déclarations :} \texttt{let x = expr;}
  \item \textbf{Affectations :} \texttt{x = expr;}
  \item \textbf{Boucles :} \texttt{repeat expr \{ ... \}}, \texttt{while (expr) \{ ... \}}
  \item \textbf{Conditions :} \texttt{if (expr) \{ ... \} else \{ ... \}}
  \item \textbf{Commandes :} \texttt{forward}, \texttt{turn}, \texttt{color}, \texttt{penDown}, \texttt{penUp}
  \item \textbf{Expressions :} littéraux, identifiants, opérateurs arithmétiques et relationnels.
\end{itemize}

\subsection*{4.2 AST (principaux nœuds)}
\begin{itemize}[leftmargin=1.2cm]
  \item \texttt{Program(body)}
  \item \texttt{VarDecl(name, value)}
  \item \texttt{Assignment(name, value)}
  \item \texttt{Repeat(count, body)}
  \item \texttt{If(condition, thenBody, elseBody)}
  \item \texttt{While(condition, body)}
  \item \texttt{Command(name, args)}
  \item \texttt{BinaryExpr(left, operator, right)}
  \item \texttt{Literal(value)}
  \item \texttt{Identifier(name)}
\end{itemize}

Chaque nœud contient une position \texttt{loc(line, col)} pour des erreurs précises.

\section*{5. Analyses : lexicale, syntaxique et sémantique}

\subsection*{5.1 Analyse lexicale}
Le lexeur transforme le texte source en une suite de \textbf{tokens}.  
Il ignore les espaces et commentaires, puis identifie : mots-clés, identifiants, nombres, chaînes, opérateurs, ponctuation.  
Chaque token porte une position \texttt{line/col} pour des erreurs précises.

\subsection*{5.2 Analyse syntaxique}
Le parseur vérifie que les tokens suivent la grammaire du langage.  
Il construit l’AST en respectant les priorités des opérateurs.  
En cas d’erreur, il tente de se resynchroniser pour continuer l’analyse.

\subsection*{5.3 Analyse sémantique}
Avant l’exécution, une passe \texttt{SemanticAnalyzer} parcourt l’AST.  
Elle vérifie les règles du langage :
\begin{itemize}[leftmargin=1.2cm]
  \item variables utilisées avant déclaration,
  \item nombre d’arguments des commandes,
  \item division par zéro quand c’est détectable statiquement.
\end{itemize}
Cette étape collecte toutes les erreurs possibles sans s’arrêter à la première.

\section*{6. Interpréteur : Logique d’exécution}
L’interpréteur visite l’AST et met à jour l’état du robot :
\begin{itemize}[leftmargin=1.2cm]
  \item variables stockées dans une table,
  \item commandes exécutées directement (mouvement, rotation, couleur),
  \item état robot : position, angle, stylo, couleur.
\end{itemize}

\section*{7. Tests et validation}
Les tests couvrent :
\begin{itemize}[leftmargin=1.2cm]
  \item \textbf{Lexeur} : types de tokens.
  \item \textbf{Parseur} : structure AST + erreurs.
  \item \textbf{Sémantique} : cas valides et invalides.
  \item \textbf{End-to-end} : pipeline complet (lexer → parser → sémantique → interprétation).
\end{itemize}

\section*{9. Défis rencontrés et solutions apportées}
\begin{itemize}[leftmargin=1.2cm]
  \item \textbf{Localisation des erreurs sémantiques :}
  \newline Ajout de \texttt{loc} dans l’AST et propagation depuis le parseur.
  \item \textbf{Analyse sémantique incomplète :}
  \newline Ajout d’une passe statique qui collecte toutes les erreurs.
  \item \textbf{Uniformité des messages :}
  \newline Harmonisation des messages et des positions.
\end{itemize}

\end{document}